\chapter{Deployment}

\section{Ambiente di sviluppo e avvio}

Il repository include un \texttt{docker-compose.dev.yml} pensato per l'esecuzione locale dell'intero sistema con i servizi effettivamente implementati. Il file di sviluppo avvia:

\begin{itemize}
\item Redis;
\item un database Mongo dedicato per auth, catalog, cart, CAD, notifications, chatbot, orders, payment;
\item i microservizi backend;
\item il frontend Vite.
\end{itemize}

La configurazione di sviluppo e' utile anche per il debug manuale di MongoDB Compass e per la verifica dei canali realtime.

\section{Composizione dei container}

Ogni servizio backend ha un proprio Dockerfile e viene eseguito in container separato. La composizione reale evita di condividere il database tra domini distinti e rende chiara la topologia dell'applicazione. Il gateway funge da facade verso il frontend, mentre il frontend dialoga quasi esclusivamente con il prefisso \texttt{/api}.

\section{Avvio locale}

Nel progetto sono presenti comandi di workspace per eseguire test backend, test e2e e build frontend. Il file root \texttt{package.json} definisce inoltre script separati per ciascun servizio. In pratica l'ambiente locale supporta tre scenari:

\begin{itemize}
\item sviluppo frontend con Vite;
\item esecuzione e test dei singoli microservizi;
\item esecuzione coordinata dell'intero stack con Docker Compose.
\end{itemize}

\section{Kubernetes (Minikube)}

Oltre a Docker Compose, il repository include manifest Kubernetes in \texttt{k8s/} per un deployment production-oriented su Minikube (namespace dedicato, immagini costruite localmente, probe di readiness/liveness). Ogni microservizio ha un manifest numerato (\texttt{20..27-*.yaml}); \texttt{payment-service} usa \texttt{26-payment-service.yaml} con Mongo dedicato (\texttt{mongo-payment}, database \texttt{paymentdb}) definito in \texttt{11-mongo-databases.yaml}.

\section{Considerazioni operative}

Il deployment attuale e' pensato per uno sviluppo riproducibile, non per una infrastruttura cloud complessa. La scelta e' coerente con lo scopo del progetto universitario: privilegiare chiarezza dell'architettura, verificabilita' dei flussi e facilita' di esecuzione locale.
